\hypertarget{documentauxe7uxe3o-da-api--cuphead-boss-guide}{%
\subsection{Documentação da API --- Cuphead Boss
Guide}\label{documentauxe7uxe3o-da-api--cuphead-boss-guide}}

API REST em JSON servida em \texttt{http://localhost:3000/api}.

\hypertarget{convenuxe7uxf5es}{%
\subsubsection{Convenções}\label{convenuxe7uxf5es}}

\hypertarget{autenticauxe7uxe3o-e-sessuxe3o}{%
\paragraph{Autenticação e sessão}\label{autenticauxe7uxe3o-e-sessuxe3o}}

\begin{itemize}
\tightlist
\item
  A autenticação é por \textbf{sessão} (cookie \texttt{cuphead.sid},
  \texttt{HttpOnly}, \texttt{SameSite=Strict}).
\item
  O navegador envia o cookie automaticamente; ferramentas como
  \texttt{curl} precisam guardá-lo
  (\texttt{-c\ cookies.txt\ -b\ cookies.txt}).
\end{itemize}

\hypertarget{proteuxe7uxe3o-csrf}{%
\paragraph{Proteção CSRF}\label{proteuxe7uxe3o-csrf}}

Toda requisição \texttt{POST}, \texttt{PUT}, \texttt{PATCH} ou
\texttt{DELETE} precisa do cabeçalho \texttt{X-CSRF-Token} com o token
obtido em \texttt{GET\ /api/auth/sessao}. O login devolve um
\textbf{token novo}, que deve substituir o anterior.

\hypertarget{formato-das-respostas}{%
\paragraph{Formato das respostas}\label{formato-das-respostas}}

Sucesso:

\begin{Verbatim}[breaklines,breakanywhere,fontsize=\scriptsize]
{ "dados": { ... } }
\end{Verbatim}

Listagens paginadas incluem \texttt{paginacao}:

\begin{Verbatim}[breaklines,breakanywhere,fontsize=\scriptsize]
{ "dados": [ ... ], "paginacao": { "pagina": 1, "limite": 6, "total": 19, "totalPaginas": 4 } }
\end{Verbatim}

Erro (todas as rotas):

\begin{Verbatim}[breaklines,breakanywhere,fontsize=\scriptsize]
{
  "erro": {
    "status": 422,
    "codigo": "VALIDACAO",
    "mensagem": "Os dados enviados são inválidos.",
    "requisicaoId": "2df0b9ae-bcc7-4214-8fa4-eb84460f3315",
    "detalhes": [{ "campo": "dificuldade", "mensagem": "deve ser menor ou igual a 5." }]
  }
}
\end{Verbatim}

Com \texttt{NODE\_ENV=development}, erros inesperados (500) também
trazem o bloco \texttt{tecnico} (nome, mensagem original e stack trace).
Em produção esse bloco nunca aparece.

\hypertarget{cuxf3digos-de-status-usados}{%
\paragraph{Códigos de status usados}\label{cuxf3digos-de-status-usados}}

\begin{longtable}[]{@{}llp{0.5\linewidth}@{}}
\toprule\noalign{}
Status & Código (\texttt{erro.codigo}) & Quando \\
\midrule\noalign{}
\endhead
\bottomrule\noalign{}
\endlastfoot
200 & --- & Leitura ou atualização bem-sucedida \\
201 & --- & Recurso criado (cabeçalho \texttt{Location} aponta para
ele) \\
204 & --- & Sucesso sem corpo (logout, exclusão, marcar conquista) \\
400 & \texttt{JSON\_INVALIDO} & Corpo da requisição não é JSON válido \\
401 & \texttt{NAO\_AUTENTICADO} & Rota exige login \\
401 & \texttt{CREDENCIAIS\_INVALIDAS} & E-mail ou senha incorretos \\
403 & \texttt{ACESSO\_NEGADO} & Perfil sem permissão (ver matriz em
\url{ARQUITETURA.md}) \\
403 & \texttt{CSRF\_INVALIDO} & Cabeçalho \texttt{X-CSRF-Token} ausente
ou incorreto \\
404 & \texttt{NAO\_ENCONTRADO} / \texttt{ROTA\_NAO\_ENCONTRADA} &
Registro ou rota inexistente \\
409 & \texttt{CONFLITO} & E-mail/nome duplicado, tática já moderada,
alterar o próprio perfil \\
413 & \texttt{CORPO\_MUITO\_GRANDE} & Corpo maior que 10 KB \\
422 & \texttt{VALIDACAO} & Tipo, tamanho ou formato inválido (lista em
\texttt{detalhes}) \\
429 & \texttt{MUITAS\_REQUISICOES} & Rate limit excedido (login ou
API) \\
500 & \texttt{ERRO\_INTERNO} & Falha inesperada (registrada no log com
\texttt{requisicaoId}) \\
\end{longtable}

\hypertarget{perfis}{%
\paragraph{Perfis}\label{perfis}}

\texttt{visitante} (sem login), \texttt{jogador} (padrão no cadastro) e
\texttt{admin}.

\begin{center}\rule{0.5\linewidth}{0.5pt}\end{center}

\hypertarget{uxedndice}{%
\subsubsection{Índice}\label{uxedndice}}

\begin{longtable}[]{@{}llp{0.5\linewidth}@{}}
\toprule\noalign{}
Método & URL & Acesso \\
\midrule\noalign{}
\endhead
\bottomrule\noalign{}
\endlastfoot
GET & \protect\hyperlink{get-apisaude}{/api/saude} & público \\
GET & \protect\hyperlink{get-apiauthsessao}{/api/auth/sessao} &
público \\
POST & \protect\hyperlink{post-apiauthregistro}{/api/auth/registro} &
público \\
POST & \protect\hyperlink{post-apiauthlogin}{/api/auth/login} & público
(rate limit) \\
POST & \protect\hyperlink{post-apiauthlogout}{/api/auth/logout} &
logado \\
GET & \protect\hyperlink{get-apiauthme}{/api/auth/me} & logado \\
GET & \protect\hyperlink{get-apiilhas}{/api/ilhas} & público \\
GET & \protect\hyperlink{get-apichefes}{/api/chefes} & público \\
GET & \protect\hyperlink{get-apichefesid}{/api/chefes/:id} & público \\
POST & \protect\hyperlink{post-apichefes}{/api/chefes} & admin \\
PUT & \protect\hyperlink{put-apichefesid}{/api/chefes/:id} & admin \\
DELETE & \protect\hyperlink{delete-apichefesid}{/api/chefes/:id} &
admin \\
GET &
\protect\hyperlink{get-apichefesidtaticas}{/api/chefes/:id/taticas} &
público \\
POST & \protect\hyperlink{post-apitaticas}{/api/taticas} & jogador,
admin \\
GET & \protect\hyperlink{get-apitaticasminhas}{/api/taticas/minhas} &
jogador, admin \\
PATCH & \protect\hyperlink{patch-apitaticasid}{/api/taticas/:id} & autor
(tática pendente) \\
DELETE & \protect\hyperlink{delete-apitaticasid}{/api/taticas/:id} &
autor (pendente) ou admin \\
GET & \protect\hyperlink{get-apitaticas}{/api/taticas} & admin \\
PATCH &
\protect\hyperlink{patch-apitaticasidstatus}{/api/taticas/:id/status} &
admin \\
GET & \protect\hyperlink{get-apiconquistas}{/api/conquistas} &
público \\
GET &
\protect\hyperlink{get-apiconquistasprogresso}{/api/conquistas/progresso}
& jogador, admin \\
PUT &
\protect\hyperlink{put-apiconquistasidprogresso}{/api/conquistas/:id/progresso}
& jogador, admin \\
DELETE &
\protect\hyperlink{delete-apiconquistasidprogresso}{/api/conquistas/:id/progresso}
& jogador, admin \\
GET & \protect\hyperlink{get-apiusuarios}{/api/usuarios} & admin \\
PATCH &
\protect\hyperlink{patch-apiusuariosidperfil}{/api/usuarios/:id/perfil}
& admin \\
\end{longtable}

\begin{center}\rule{0.5\linewidth}{0.5pt}\end{center}

\hypertarget{sauxfade}{%
\subsubsection{Saúde}\label{sauxfade}}

\hypertarget{get-apisaude}{%
\paragraph{GET /api/saude}\label{get-apisaude}}

Verifica se a API e o banco estão respondendo.

\textbf{Parâmetros:} nenhum.

\begin{Verbatim}[breaklines,breakanywhere,fontsize=\scriptsize]
curl http://localhost:3000/api/saude
\end{Verbatim}

\begin{Verbatim}[breaklines,breakanywhere,fontsize=\scriptsize]
{ "dados": { "status": "ok", "banco": "ok" } }
\end{Verbatim}

\textbf{Status:} \texttt{200}, \texttt{500} (banco indisponível).

\begin{center}\rule{0.5\linewidth}{0.5pt}\end{center}

\hypertarget{autenticauxe7uxe3o}{%
\subsubsection{Autenticação}\label{autenticauxe7uxe3o}}

\hypertarget{get-apiauthsessao}{%
\paragraph{GET /api/auth/sessao}\label{get-apiauthsessao}}

Retorna o usuário logado (ou \texttt{null}) e o token CSRF da sessão. É
a primeira chamada do frontend.

\textbf{Parâmetros:} nenhum.

\begin{Verbatim}[breaklines,breakanywhere,fontsize=\scriptsize]
curl -c cookies.txt -b cookies.txt http://localhost:3000/api/auth/sessao
\end{Verbatim}

\begin{Verbatim}[breaklines,breakanywhere,fontsize=\scriptsize]
{ "dados": { "usuario": null, "csrfToken": "I8bheqSPZxh4GORdVYSbqZJvIHBMyNJhKqCkg-5M5AM" } }
\end{Verbatim}

\textbf{Status:} \texttt{200}.

\hypertarget{post-apiauthregistro}{%
\paragraph{POST /api/auth/registro}\label{post-apiauthregistro}}

Cria uma conta com perfil \texttt{jogador}. Campos extras (ex.:
\texttt{perfil}) são ignorados.

\textbf{Corpo (JSON):}

\begin{longtable}[]{@{}llp{0.5\linewidth}@{}}
\toprule\noalign{}
Campo & Tipo & Regras \\
\midrule\noalign{}
\endhead
\bottomrule\noalign{}
\endlastfoot
\texttt{nome} & texto & obrigatório, 2--60 caracteres \\
\texttt{email} & texto & obrigatório, e-mail válido, até 254 caracteres
(convertido para minúsculas) \\
\texttt{senha} & texto & obrigatório, 8+ caracteres, letras e números,
até 72 bytes \\
\end{longtable}

\begin{Verbatim}[breaklines,breakanywhere,fontsize=\scriptsize]
curl -b cookies.txt -c cookies.txt -X POST http://localhost:3000/api/auth/registro \
  -H "Content-Type: application/json" -H "X-CSRF-Token: <token>" \
  -d '{"nome":"Cuphead","email":"cup@head.com","senha":"xicara123"}'
\end{Verbatim}

\begin{Verbatim}[breaklines,breakanywhere,fontsize=\scriptsize]
{
  "dados": { "id": 2, "nome": "Cuphead", "email": "cup@head.com", "perfil": "jogador", "criadoEm": "2026-10-14T17:13:40.620Z" }
}
\end{Verbatim}

\textbf{Status:} \texttt{201}, \texttt{403} (CSRF), \texttt{409} (e-mail
já cadastrado), \texttt{422}.

\hypertarget{post-apiauthlogin}{%
\paragraph{POST /api/auth/login}\label{post-apiauthlogin}}

Autentica e cria a sessão. O id de sessão é regenerado (anti-fixação) e
um \textbf{novo token CSRF} é devolvido. Limitado a
\texttt{RATE\_LIMIT\_LOGIN\_MAX} tentativas malsucedidas por IP a cada
\texttt{RATE\_LIMIT\_LOGIN\_JANELA\_MIN} minutos.

\textbf{Corpo (JSON):} \texttt{email} (texto, obrigatório),
\texttt{senha} (texto, obrigatório).

\begin{Verbatim}[breaklines,breakanywhere,fontsize=\scriptsize]
curl -b cookies.txt -c cookies.txt -X POST http://localhost:3000/api/auth/login \
  -H "Content-Type: application/json" -H "X-CSRF-Token: <token>" \
  -d '{"email":"cup@head.com","senha":"xicara123"}'
\end{Verbatim}

\begin{Verbatim}[breaklines,breakanywhere,fontsize=\scriptsize]
{
  "dados": {
    "usuario": { "id": 2, "nome": "Cuphead", "email": "cup@head.com", "perfil": "jogador", "criadoEm": "2026-10-14T17:13:40.620Z" },
    "csrfToken": "RMBNkIqvIVZ8zJvx7kQvLUQUd6Mnrq8bDTWoCPp2sp4"
  }
}
\end{Verbatim}

Credenciais erradas (mesma resposta para e-mail inexistente ou senha
incorreta):

\begin{Verbatim}[breaklines,breakanywhere,fontsize=\scriptsize]
{ "erro": { "status": 401, "codigo": "CREDENCIAIS_INVALIDAS", "mensagem": "E-mail ou senha incorretos.", "requisicaoId": "f387e924-299a-423b-a5b2-50c3b2ecc659" } }
\end{Verbatim}

\textbf{Status:} \texttt{200}, \texttt{401}, \texttt{403} (CSRF),
\texttt{422}, \texttt{429}.

\hypertarget{post-apiauthlogout}{%
\paragraph{POST /api/auth/logout}\label{post-apiauthlogout}}

Destroi a sessão no servidor e apaga o cookie. O cookie antigo deixa de
valer.

\textbf{Parâmetros:} nenhum (exige \texttt{X-CSRF-Token}).

\begin{Verbatim}[breaklines,breakanywhere,fontsize=\scriptsize]
curl -b cookies.txt -c cookies.txt -X POST http://localhost:3000/api/auth/logout -H "X-CSRF-Token: <token>"
\end{Verbatim}

Resposta: corpo vazio.

\textbf{Status:} \texttt{204}, \texttt{401}, \texttt{403} (CSRF).

\hypertarget{get-apiauthme}{%
\paragraph{GET /api/auth/me}\label{get-apiauthme}}

Dados do usuário logado.

\begin{Verbatim}[breaklines,breakanywhere,fontsize=\scriptsize]
curl -b cookies.txt http://localhost:3000/api/auth/me
\end{Verbatim}

\begin{Verbatim}[breaklines,breakanywhere,fontsize=\scriptsize]
{ "dados": { "id": 2, "nome": "Cuphead", "email": "cup@head.com", "perfil": "jogador", "criadoEm": "2026-10-14T17:13:40.620Z" } }
\end{Verbatim}

\textbf{Status:} \texttt{200}, \texttt{401}.

\begin{center}\rule{0.5\linewidth}{0.5pt}\end{center}

\hypertarget{ilhas}{%
\subsubsection{Ilhas}\label{ilhas}}

\hypertarget{get-apiilhas}{%
\paragraph{GET /api/ilhas}\label{get-apiilhas}}

\begin{Verbatim}[breaklines,breakanywhere,fontsize=\scriptsize]
curl http://localhost:3000/api/ilhas
\end{Verbatim}

\begin{Verbatim}[breaklines,breakanywhere,fontsize=\scriptsize]
{
  "dados": [
    { "id": 1, "nome": "Ilha 1", "ordem": 1 },
    { "id": 2, "nome": "Ilha 2", "ordem": 2 },
    { "id": 3, "nome": "Ilha 3", "ordem": 3 },
    { "id": 4, "nome": "Inkwell Hell", "ordem": 4 }
  ]
}
\end{Verbatim}

\textbf{Status:} \texttt{200}.

\begin{center}\rule{0.5\linewidth}{0.5pt}\end{center}

\hypertarget{chefes}{%
\subsubsection{Chefes}\label{chefes}}

Objeto \texttt{Chefe}:

\begin{Verbatim}[breaklines,breakanywhere,fontsize=\scriptsize]
{
  "id": 1,
  "nome": "The Root Pack",
  "fase": "Botanic Panic",
  "estrategia": "O FSM transita de forma linear: Batata (projéteis horizontais), ...",
  "dificuldade": 1,
  "ilhaId": 1,
  "ilhaNome": "Ilha 1",
  "criadoEm": "2026-10-14T17:12:36.512Z",
  "atualizadoEm": "2026-10-14T17:12:36.512Z",
  "totalTaticas": 1
}
\end{Verbatim}

\hypertarget{get-apichefes}{%
\paragraph{GET /api/chefes}\label{get-apichefes}}

Lista paginada, ordenada por ilha.

\textbf{Query string:}

\begin{longtable}[]{@{}lllp{0.42\linewidth}@{}}
\toprule\noalign{}
Parâmetro & Tipo & Padrão & Regras \\
\midrule\noalign{}
\endhead
\bottomrule\noalign{}
\endlastfoot
\texttt{pagina} & inteiro & 1 & ≥ 1 \\
\texttt{limite} & inteiro & 6 & 1--50 \\
\texttt{ilha} & inteiro & --- & id da ilha \\
\texttt{busca} & texto & --- & 1--60 caracteres; procura em nome e fase,
sem diferenciar maiúsculas \\
\end{longtable}

\begin{Verbatim}[breaklines,breakanywhere,fontsize=\scriptsize]
curl "http://localhost:3000/api/chefes?ilha=4&busca=dev"
\end{Verbatim}

\begin{Verbatim}[breaklines,breakanywhere,fontsize=\scriptsize]
{
  "dados": [
    { "id": 19, "nome": "The Devil", "fase": "One Hell of a Time", "estrategia": "Fase 1: Leia os olhos dele ...", "dificuldade": 5,
      "ilhaId": 4, "ilhaNome": "Inkwell Hell", "criadoEm": "...", "atualizadoEm": "...", "totalTaticas": 0 }
  ],
  "paginacao": { "pagina": 1, "limite": 6, "total": 1, "totalPaginas": 1 }
}
\end{Verbatim}

\textbf{Status:} \texttt{200}, \texttt{422} (ex.: \texttt{limite=999}).

\hypertarget{get-apichefesid}{%
\paragraph{GET /api/chefes/:id}\label{get-apichefesid}}

\textbf{Parâmetros de rota:} \texttt{id} (inteiro ≥ 1).

\begin{Verbatim}[breaklines,breakanywhere,fontsize=\scriptsize]
curl http://localhost:3000/api/chefes/1
\end{Verbatim}

Resposta: \texttt{\{\ "dados":\ Chefe\ \}}.

\textbf{Status:} \texttt{200}, \texttt{404}, \texttt{422} (id não
numérico).

\hypertarget{post-apichefes}{%
\paragraph{POST /api/chefes}\label{post-apichefes}}

Somente \textbf{admin}.

\textbf{Corpo (JSON):}

\begin{longtable}[]{@{}llp{0.5\linewidth}@{}}
\toprule\noalign{}
Campo & Tipo & Regras \\
\midrule\noalign{}
\endhead
\bottomrule\noalign{}
\endlastfoot
\texttt{nome} & texto & obrigatório, 2--60, único (sem diferenciar
maiúsculas) \\
\texttt{fase} & texto & obrigatório, 2--80 \\
\texttt{ilhaId} & inteiro & obrigatório, ilha existente \\
\texttt{estrategia} & texto & obrigatório, 20--1000, aceita quebras de
linha \\
\texttt{dificuldade} & inteiro & obrigatório, 1--5 \\
\end{longtable}

\begin{Verbatim}[breaklines,breakanywhere,fontsize=\scriptsize]
curl -b cookies.txt -X POST http://localhost:3000/api/chefes \
  -H "Content-Type: application/json" -H "X-CSRF-Token: <token>" \
  -d '{"nome":"Glumstone the Giant","fase":"Doggone Dogfight","ilhaId":4,"estrategia":"Atire nos olhos e evite as mãos.","dificuldade":4}'
\end{Verbatim}

\begin{Verbatim}[breaklines,breakanywhere,fontsize=\scriptsize]
{
  "dados": { "id": 20, "nome": "Glumstone the Giant", "fase": "Doggone Dogfight", "estrategia": "Atire nos olhos e evite as mãos.",
             "dificuldade": 4, "ilhaId": 4, "ilhaNome": "Inkwell Hell", "criadoEm": "...", "atualizadoEm": "...", "totalTaticas": 0 }
}
\end{Verbatim}

Cabeçalho \texttt{Location:\ /api/chefes/20}.

\textbf{Status:} \texttt{201}, \texttt{401}, \texttt{403}, \texttt{409}
(nome duplicado), \texttt{422}.

\hypertarget{put-apichefesid}{%
\paragraph{PUT /api/chefes/:id}\label{put-apichefesid}}

Somente \textbf{admin}. Substitui todos os campos (mesmo corpo do
\texttt{POST}).

\begin{Verbatim}[breaklines,breakanywhere,fontsize=\scriptsize]
curl -b cookies.txt -X PUT http://localhost:3000/api/chefes/20 \
  -H "Content-Type: application/json" -H "X-CSRF-Token: <token>" \
  -d '{"nome":"Glumstone","fase":"Doggone Dogfight","ilhaId":4,"estrategia":"Estratégia atualizada pelo administrador.","dificuldade":3}'
\end{Verbatim}

Resposta: \texttt{\{\ "dados":\ Chefe\ \}} atualizado.

\textbf{Status:} \texttt{200}, \texttt{401}, \texttt{403}, \texttt{404},
\texttt{409}, \texttt{422}.

\hypertarget{delete-apichefesid}{%
\paragraph{DELETE /api/chefes/:id}\label{delete-apichefesid}}

Somente \textbf{admin}. Remove o chefe e, em cascata, as táticas dele.

\begin{Verbatim}[breaklines,breakanywhere,fontsize=\scriptsize]
curl -b cookies.txt -X DELETE http://localhost:3000/api/chefes/20 -H "X-CSRF-Token: <token>"
\end{Verbatim}

Resposta: corpo vazio.

\textbf{Status:} \texttt{204}, \texttt{401}, \texttt{403}, \texttt{404}.

\hypertarget{get-apichefesidtaticas}{%
\paragraph{GET /api/chefes/:id/taticas}\label{get-apichefesidtaticas}}

Táticas \textbf{aprovadas} do chefe, mais recentes primeiro.

\begin{Verbatim}[breaklines,breakanywhere,fontsize=\scriptsize]
curl http://localhost:3000/api/chefes/1/taticas
\end{Verbatim}

\begin{Verbatim}[breaklines,breakanywhere,fontsize=\scriptsize]
{
  "dados": [
    { "id": 1, "conteudo": "Fique no canto durante a fase da cenoura.", "status": "aprovada", "chefeId": 1, "chefeNome": "The Root Pack",
      "usuarioId": 2, "autorNome": "Cuphead", "criadoEm": "...", "atualizadoEm": "..." }
  ]
}
\end{Verbatim}

\textbf{Status:} \texttt{200}, \texttt{404}, \texttt{422}.

\begin{center}\rule{0.5\linewidth}{0.5pt}\end{center}

\hypertarget{tuxe1ticas-da-comunidade}{%
\subsubsection{Táticas da comunidade}\label{tuxe1ticas-da-comunidade}}

Objeto \texttt{Tatica}: \texttt{id}, \texttt{conteudo}, \texttt{status}
(\texttt{pendente} \textbar{} \texttt{aprovada} \textbar{}
\texttt{rejeitada}), \texttt{chefeId}, \texttt{chefeNome},
\texttt{usuarioId}, \texttt{autorNome}, \texttt{criadoEm},
\texttt{atualizadoEm}.

\hypertarget{post-apitaticas}{%
\paragraph{POST /api/taticas}\label{post-apitaticas}}

Jogador ou admin envia uma tática. Ela nasce \texttt{pendente}.

\textbf{Corpo (JSON):} \texttt{chefeId} (inteiro, chefe existente),
\texttt{conteudo} (texto, 10--1000, aceita quebras de linha).

\begin{Verbatim}[breaklines,breakanywhere,fontsize=\scriptsize]
curl -b cookies.txt -X POST http://localhost:3000/api/taticas \
  -H "Content-Type: application/json" -H "X-CSRF-Token: <token>" \
  -d '{"chefeId":1,"conteudo":"Fique no canto durante a fase da cenoura."}'
\end{Verbatim}

\begin{Verbatim}[breaklines,breakanywhere,fontsize=\scriptsize]
{
  "dados": { "id": 1, "conteudo": "Fique no canto durante a fase da cenoura.", "status": "pendente", "chefeId": 1,
             "chefeNome": "The Root Pack", "usuarioId": 2, "autorNome": "Cuphead", "criadoEm": "...", "atualizadoEm": "..." }
}
\end{Verbatim}

\textbf{Status:} \texttt{201}, \texttt{401}, \texttt{403} (CSRF),
\texttt{422} (inclui chefe inexistente).

\hypertarget{get-apitaticasminhas}{%
\paragraph{GET /api/taticas/minhas}\label{get-apitaticasminhas}}

Todas as táticas do usuário logado (qualquer status).

\begin{Verbatim}[breaklines,breakanywhere,fontsize=\scriptsize]
curl -b cookies.txt http://localhost:3000/api/taticas/minhas
\end{Verbatim}

Resposta: \texttt{\{\ "dados":\ {[}Tatica,\ ...{]}\ \}}.

\textbf{Status:} \texttt{200}, \texttt{401}.

\hypertarget{patch-apitaticasid}{%
\paragraph{PATCH /api/taticas/:id}\label{patch-apitaticasid}}

O \textbf{autor} altera o conteúdo enquanto a tática está
\texttt{pendente}.

\textbf{Corpo (JSON):} \texttt{conteudo} (texto, 10--1000).

\begin{Verbatim}[breaklines,breakanywhere,fontsize=\scriptsize]
curl -b cookies.txt -X PATCH http://localhost:3000/api/taticas/1 \
  -H "Content-Type: application/json" -H "X-CSRF-Token: <token>" \
  -d '{"conteudo":"Conteúdo editado pelo autor enquanto pendente."}'
\end{Verbatim}

Resposta: \texttt{\{\ "dados":\ Tatica\ \}}.

\textbf{Status:} \texttt{200}, \texttt{401}, \texttt{403} (não é o
autor), \texttt{404}, \texttt{409} (já moderada), \texttt{422}.

\hypertarget{delete-apitaticasid}{%
\paragraph{DELETE /api/taticas/:id}\label{delete-apitaticasid}}

O autor exclui a própria tática pendente; o \textbf{admin} exclui
qualquer uma.

\begin{Verbatim}[breaklines,breakanywhere,fontsize=\scriptsize]
curl -b cookies.txt -X DELETE http://localhost:3000/api/taticas/1 -H "X-CSRF-Token: <token>"
\end{Verbatim}

\textbf{Status:} \texttt{204}, \texttt{401}, \texttt{403}, \texttt{404},
\texttt{409}.

\hypertarget{get-apitaticas}{%
\paragraph{GET /api/taticas}\label{get-apitaticas}}

Fila de moderação (somente \textbf{admin}), mais antigas primeiro.

\textbf{Query string:} \texttt{status} (\texttt{pendente} \textbar{}
\texttt{aprovada} \textbar{} \texttt{rejeitada}, opcional),
\texttt{pagina}, \texttt{limite} (padrão 20, máx. 50).

\begin{Verbatim}[breaklines,breakanywhere,fontsize=\scriptsize]
curl -b cookies.txt "http://localhost:3000/api/taticas?status=pendente"
\end{Verbatim}

Resposta:
\texttt{\{\ "dados":\ {[}Tatica,\ ...{]},\ "paginacao":\ \{\ ...\ \}\ \}}.

\textbf{Status:} \texttt{200}, \texttt{401}, \texttt{403}, \texttt{422}.

\hypertarget{patch-apitaticasidstatus}{%
\paragraph{PATCH
/api/taticas/:id/status}\label{patch-apitaticasidstatus}}

Aprova ou rejeita uma tática (somente \textbf{admin}).

\textbf{Corpo (JSON):} \texttt{status} (\texttt{aprovada} \textbar{}
\texttt{rejeitada}).

\begin{Verbatim}[breaklines,breakanywhere,fontsize=\scriptsize]
curl -b cookies.txt -X PATCH http://localhost:3000/api/taticas/1/status \
  -H "Content-Type: application/json" -H "X-CSRF-Token: <token>" -d '{"status":"aprovada"}'
\end{Verbatim}

Resposta: \texttt{\{\ "dados":\ Tatica\ \}} com o novo status.

\textbf{Status:} \texttt{200}, \texttt{401}, \texttt{403}, \texttt{404},
\texttt{422}.

\begin{center}\rule{0.5\linewidth}{0.5pt}\end{center}

\hypertarget{conquistas}{%
\subsubsection{Conquistas}\label{conquistas}}

\hypertarget{get-apiconquistas}{%
\paragraph{GET /api/conquistas}\label{get-apiconquistas}}

Catálogo com as 28 conquistas.

\begin{Verbatim}[breaklines,breakanywhere,fontsize=\scriptsize]
curl http://localhost:3000/api/conquistas
\end{Verbatim}

\begin{Verbatim}[breaklines,breakanywhere,fontsize=\scriptsize]
{ "dados": [ { "id": 1, "nome": "Taking Names", "descricao": "Derrote um chefão." }, "..." ] }
\end{Verbatim}

\textbf{Status:} \texttt{200}.

\hypertarget{get-apiconquistasprogresso}{%
\paragraph{GET
/api/conquistas/progresso}\label{get-apiconquistasprogresso}}

Progresso do usuário logado.

\begin{Verbatim}[breaklines,breakanywhere,fontsize=\scriptsize]
curl -b cookies.txt http://localhost:3000/api/conquistas/progresso
\end{Verbatim}

\begin{Verbatim}[breaklines,breakanywhere,fontsize=\scriptsize]
{ "dados": { "total": 28, "concluidas": 1, "percentual": 4, "itens": [ { "conquistaId": 3, "concluidaEm": "2026-10-14T17:13:41.244Z" } ] } }
\end{Verbatim}

\textbf{Status:} \texttt{200}, \texttt{401}.

\hypertarget{put-apiconquistasidprogresso}{%
\paragraph{PUT
/api/conquistas/:id/progresso}\label{put-apiconquistasidprogresso}}

Marca a conquista como concluída (idempotente: repetir não gera erro).

\begin{Verbatim}[breaklines,breakanywhere,fontsize=\scriptsize]
curl -b cookies.txt -X PUT http://localhost:3000/api/conquistas/3/progresso -H "X-CSRF-Token: <token>"
\end{Verbatim}

\textbf{Status:} \texttt{204}, \texttt{401}, \texttt{403} (CSRF),
\texttt{404}, \texttt{422}.

\hypertarget{delete-apiconquistasidprogresso}{%
\paragraph{DELETE
/api/conquistas/:id/progresso}\label{delete-apiconquistasidprogresso}}

Desmarca a conquista.

\begin{Verbatim}[breaklines,breakanywhere,fontsize=\scriptsize]
curl -b cookies.txt -X DELETE http://localhost:3000/api/conquistas/3/progresso -H "X-CSRF-Token: <token>"
\end{Verbatim}

\textbf{Status:} \texttt{204}, \texttt{401}, \texttt{403} (CSRF),
\texttt{404}, \texttt{422}.

\begin{center}\rule{0.5\linewidth}{0.5pt}\end{center}

\hypertarget{usuuxe1rios-administrauxe7uxe3o}{%
\subsubsection{Usuários
(administração)}\label{usuuxe1rios-administrauxe7uxe3o}}

\hypertarget{get-apiusuarios}{%
\paragraph{GET /api/usuarios}\label{get-apiusuarios}}

Lista paginada de usuários (somente \textbf{admin}). Nunca inclui hash
de senha.

\textbf{Query string:} \texttt{pagina}, \texttt{limite} (padrão 20, máx.
50).

\begin{Verbatim}[breaklines,breakanywhere,fontsize=\scriptsize]
curl -b cookies.txt http://localhost:3000/api/usuarios
\end{Verbatim}

\begin{Verbatim}[breaklines,breakanywhere,fontsize=\scriptsize]
{
  "dados": [ { "id": 1, "nome": "Administrador", "email": "admin@cuphead.local", "perfil": "admin", "criadoEm": "..." } ],
  "paginacao": { "pagina": 1, "limite": 20, "total": 1, "totalPaginas": 1 }
}
\end{Verbatim}

\textbf{Status:} \texttt{200}, \texttt{401}, \texttt{403}, \texttt{422}.

\hypertarget{patch-apiusuariosidperfil}{%
\paragraph{PATCH
/api/usuarios/:id/perfil}\label{patch-apiusuariosidperfil}}

Altera o perfil de outro usuário (somente \textbf{admin}). Vale
imediatamente, inclusive para sessões já abertas.

\textbf{Corpo (JSON):} \texttt{perfil} (\texttt{jogador} \textbar{}
\texttt{admin}).

\begin{Verbatim}[breaklines,breakanywhere,fontsize=\scriptsize]
curl -b cookies.txt -X PATCH http://localhost:3000/api/usuarios/2/perfil \
  -H "Content-Type: application/json" -H "X-CSRF-Token: <token>" -d '{"perfil":"admin"}'
\end{Verbatim}

Resposta: \texttt{\{\ "dados":\ Usuario\ \}} atualizado.

\textbf{Status:} \texttt{200}, \texttt{401}, \texttt{403}, \texttt{404},
\texttt{409} (alterar o próprio perfil), \texttt{422}.
